\begin{table}\centering\begin{threeparttable}
\caption{Descriptive statistics.}\label{tab:ds}
\begin{tabular}{lccccc} \hline \hline
& Mean & SD  & Min  & Max  \\  \hline 
\textit{\underline{Municipal level}} & & & & \\
TV tax compliance&     55.0&     28.9&        0&    100.0\\
Referendum turnout&     68.4&     21.6&     10.3&      100\\
European ele. turnout&     75.0&     14.0&     13.7&      100\\
N. civic association * 1000&     2.96&     3.42&        0&     61.2\\
Population      &   6779.3&  40873.8&       30&  2840259\\
Elderly index   &    113.8&    116.2&     5.11&     4150\\
Education gender gap index&    137.8&     77.9&        0&   2640.3\\
Illiterate index&     5.00&     7.16&        0&     63.6\\
Workforce participation&     50.0&     7.45&     16.7&     89.6\\
& & & & & \\
%
%
\textit{\underline{Provincial level}} & & & & \\
N. blood donations&    0.030&    0.022&        0&     0.11\\
N. CONI associations / 1000 people&     1.30&     0.41&     0.43&     2.45\\
N. of non-sport newspapers sold / 1000&     80.1&     35.1&     17.9&    175.4\\
Trust (WVS)     &     0.32&     0.14&        0&     0.85\\
\hline
\end{tabular} \begin{tablenotes} \item \scriptsize{\textit{\underline{Notes}.} This table shows the descriptive statistics of the main variables at our disposal in the analysis. } \end{tablenotes} \end{threeparttable} \end{table}
